%\begin{minipage}{0.80\textwidth}
%\begin{algorithm2e}[H]
%    \dontprintsemicolon
%    \linesnumbered
%    \SetLine
%    \BlankLine
%    \Entrada{$X, Y, num\_samples$}
%    \Saida{$bootstrap\_samples$}
%    \BlankLine
%    \Inicio{
%        %Definir $n \gets \text{size}(X, 1)$\;
%        %Criar uma célula $bootstrap\_samples$ de %dimensão $num\_samples$ por 2\;
%				\ParaTodo{$i \gets 1$:$num\_samples$}
%        %\Para{$i \gets 1$ \To $num\_samples$}
%				{
%            Gerar índices aleatórios: $indices \gets %\text{randi}([1, n], n, 1)$\;
%            Criar amostras:\;
%            $X\_bootstrap \gets X(indices, :)$\;
%            $Y\_bootstrap \gets Y(indices)$\;
%            Armazenar amostras:\;
%            $bootstrap\_samples(i, 1) \gets %X\_bootstrap$\;
%            $bootstrap\_samples(i, 2) \gets %Y\_bootstrap$\;
%        }
%        \Retorna{$bootstrap\_samples$}\;
%    }
%    \caption{Estimativa das replicações %\textit{bootstrap}.\label{alg:bootstrap2}}
%\end{algorithm2e}
%\end{minipage}

%----------------------------------------------
\begin{comment}

\begin{minipage}{0.80\textwidth}
\begin{algorithm2e}[H]
    \dontprintsemicolon
    \linesnumbered
    \SetLine
    \BlankLine
    \Entrada{$X, Y, m$}
    \Saida{$bootstrap\_samples$}
    \BlankLine
    \Inicio{
        Definir $n \gets \text{size}(X, 1)$\;
        Criar uma célula $bootstrap\_samples$ de dimensão $m$ por 2\;
        Selecione a quantidade de amostras independentes $m$,~\{$x^{*1}, x^{*2}, \ldots, x^{*m}$\}, tendo tamanho $m$ e podendo haver repetições de $x$\;
        \ParaTodo{$b \gets 1$:$m$}{
            Gerar índices aleatórios: $indices \gets \text{randi}([1, n], n, 1)$\;
            Criar amostras:\;
            $X\_bootstrap \gets X(indices, :)$\;
            $Y\_bootstrap \gets Y(indices)$\;
            Avaliar a replicação \textit{bootstrap}: $\widehat{\theta}^*(b) \gets s(x^{*b})$\;
            Armazenar amostras:\;
            $bootstrap\_samples(b, 1) \gets X\_bootstrap$\;
            $bootstrap\_samples(b, 2) \gets Y\_bootstrap$\;
        }
        \Retorna{$bootstrap\_samples$}\;
    }
    \caption{Estimativa das replicações \textit{bootstrap} combinada.\label{alg:bootstrap3}}
\end{algorithm2e}
\end{minipage}

\end{comment}
%-----------------------------------------------------

\begin{comment}
%--------------------------------------------------
\begin{minipage}{0.80\textwidth}
\begin{algorithm2e}[H]
 \dontprintsemicolon
 \linesnumbered
 %%\nocaptionofalgo
 \SetLine
 \BlankLine
 \Entrada{Amostra original $x = (x_1, x_2, \ldots, x_n)$, estatística de interesse $\theta = s(x)$, número de replicações bootstrap $m$}
 \Saida{Estimativa bootstrap $\widehat{\theta}^*$, erro padrão bootstrap, intervalo de confiança bootstrap}
 \BlankLine
\Inicio{
  Selecione a quantidade de amostras independentes ~$m$,~\{$x^{*1},x^{*2},\ldots,x^{*m}$\}, tendo tamanho
 $m$ e podendo haver repetições de  $x$.\;
Avaliar a replicação \textit{bootstrap} para amostra \textit{bootstrap}, $\widehat{\theta}^*(b)= s(x^{*b})$ $\quad$ $b = 1,2,\ldots ,m$. \;}
\caption{Estimativa das replicações \textit{bootstrap}. \label{alg:bootstrap}}
\end{algorithm2e}
\end{minipage}
%------------------------------------------------
\end{comment}


\begin{comment}
%----------------------------------------------------
\begin{minipage}{0.80\textwidth}
\begin{algorithm2e}[H]
 \dontprintsemicolon
 \linesnumbered
 %%\nocaptionofalgo
 \SetLine
 \BlankLine
 \Entrada{Amostra original $x = (x_1, x_2, \ldots, x_n)$, estatística de interesse $\theta = s(x)$, número de amostras bootstrap $m$, tamanho das amostras bootstrap (opcional, padrão $n$)}
 \Saida{Estimativa bagging $\widehat{\theta}^{*}_{bag}$, distribuição das estimativas bootstrap $\{\widehat{\theta}(1), \ldots, \widehat{\theta}(m)\}$}
 \BlankLine
\Inicio{
Selecione a quantidade de amostras \textit{bootstrap} ~$m$,~\{$x^{*1},x^{*2},\ldots,x^{*m}$\}, podendo haver repetições de  $x$.\;
Selecione o tamanho $n$ que cada amostra \textit{bootstrap} terá, no caso do \textit{bagging} tradicional, todas amostras \textit{bootstrap} tem o mesmo tamanho que a amostra original.\;
\Repita{$(b=m)$\;}{Avaliar a replicação \textit{bootstrap }para cada amostra \textit{bootstrap},\;
$\widehat{\theta}(b)= s(x^{*b}) ~~b = 1,2,\ldots,m.$\;}
Agregar as replicações para amostra \textit{bootstrap} $\widehat{\theta}(b)$ em $\widehat{\theta}^{*}_{bag}$,\;
$\widehat{\theta}^{*}_{bag} = \frac{1}{n}\sum_{i=1}^{n}s(x^{*b})~~b = 1,2,\ldots,m.$\; 
}
\caption{A Estimativa \textit{bagging}. \label{alg:bagging}}
\end{algorithm2e}
\end{minipage}
%---------------------------------------------------
\end{comment}

\begin{comment}

\begin{minipage}{0.80\textwidth}
\scalebox{0.85}{
\begin{algorithm2e}[H]
 \dontprintsemicolon
 \linesnumbered
 \SetLine
 \BlankLine
 \Entrada{Amostra original $x = (x_1, x_2, \ldots, x_n)$, estatística de interesse $\theta = s(x)$, \\
 \hspace{1.5em}número de amostras bootstrap $m$, \\
 \hspace{1.5em}tamanho máximo das subamostras $N_{\text{Bootstrap}}$, \\
 \hspace{1.5em}variação do tamanho das subamostras (opcional)}
 \Saida{Estimativa subagging $\widehat{\theta}^{*}_{subag}$, \\
 \hspace{1.5em}conjunto de estimativas parciais $\{\widehat{\theta}(1), \ldots, \widehat{\theta}(m)\}$, \\
 \hspace{1.5em}erro padrão da estimativa}
 \BlankLine
Selecione o tamanho máximo que as amostras \textit{bootstrap} \textit{N.Bootstrap} poderão ter.\;
Selecione a quantidade de amostras \textit{bootstrap} independentes $m$, $\{x^{*1}, x^{*2}, \ldots, x^{*m}\}$, não podendo haver repetições de $x$.\;
Selecione a variação do tamanho \textit{N.Bootstrap} que cada amostra \textit{bootstrap} terá, no \textit{subagging }as amostras \textit{bootstrap} não necessitam ter o mesmo tamanho que a amostra original.\;
 \BlankLine
 \Repita{(\textit{N.Bootstrap})}{\Repita{($b=m$)}{Avaliar a replicação \textit{bootstrap} para cada amostra \textit{bootstrap}, $\widehat{\theta}(b)= s(x^{*b}) ~~b = 1,2,\ldots,m.$\;}}
Agregar as replicações para cada amostra \textit{bootstrap} $\widehat{\theta}(b)$ em $\widehat{\theta}^{*}_{bag}$,\;
$\widehat{\theta}^{*}_{subag} = \dfrac{1}{n}\sum_{i=1}^{n}s(x^{*b})~~b = 1,2,\ldots,m-1.$\;
\caption{Estimativa \textit{subagging}. \label{alg:subagging}}
\end{algorithm2e}
}
\end{minipage}

\end{comment}